\documentclass[12pt,a4paper]{article}
\usepackage[utf8]{inputenc}
\usepackage[T1]{fontenc}
\usepackage[croatian]{babel}
\usepackage{geometry}
\usepackage{graphicx}
\usepackage{setspace}
\usepackage{titlesec}
\usepackage{fancyhdr}
\usepackage{caption}
\usepackage{longtable}
\usepackage{array}
\usepackage{tabularx}
\usepackage{amsmath}
\usepackage{amssymb}
\usepackage{mathptmx}       % Times New Roman font
\usepackage{chngcntr}       % Za numeriranje po poglavljima

% Postavke margina
\geometry{left=2.5cm,right=2.5cm,top=2.5cm,bottom=2.5cm}

% Prored 1.5
\onehalfspacing

% ---- STILOVI NASLOVA ----
% Section: nova stranica + bold uppercase
\titleformat{\section}
  {\normalfont\large\bfseries\uppercase}{\thesection.}{0.5em}{}
\titleformat{\subsection}
  {\normalfont\normalsize\bfseries}{\thesubsection.}{0.5em}{}
\titleformat{\subsubsection}
  {\normalfont\normalsize\bfseries\itshape}{\thesubsubsection.}{0.5em}{}
\titleformat{\paragraph}[hang]
  {\normalfont\normalsize\itshape}{\theparagraph.}{0.5em}{}
\titlespacing*{\paragraph}{0pt}{3.25ex plus 1ex minus .2ex}{0.5em}

% Svaki \section počinje na novoj stranici
\newcommand{\sectionbreak}{\clearpage}

% ---- NUMERIRANJE SLIKA, TABLICA I JEDNADŽBI PO POGLAVLJIMA ----
\counterwithin{figure}{section}
\counterwithin{table}{section}
\counterwithin{equation}{section}

% Dubina numeriranja (do 4 razine)
\setcounter{secnumdepth}{4}
\setcounter{tocdepth}{4}

% ---- NAZIVI AUTOMATSKIH POPISA (addto jer babel croatian overrida renewcommand) ----
\addto\captionscroatian{%
  \renewcommand{\contentsname}{SADRŽAJ}%
  \renewcommand{\listfigurename}{POPIS SLIKA}%
  \renewcommand{\listtablename}{POPIS TABLICA}%
  \renewcommand{\refname}{LITERATURA}%
}

% ---- POTPISI SLIKA I TABLICA ----
% Captions: 11pt Bold
\captionsetup[figure]{font={small,bf},labelsep=period}
\captionsetup[table]{font={small,bf},labelsep=period,position=above}

% ---- ZAGLAVLJE I PODNOŽJE ----
\pagestyle{fancy}
\fancyhf{}
\fancyhead[L]{\small\textit{{{IME_I_PREZIME}}}}
\fancyhead[R]{\small\textit{Seminar}}
\fancyfoot[L]{\small\textit{Fakultet strojarstva i brodogradnje}}
\fancyfoot[R]{\small\textit{\thepage}}
\renewcommand{\headrulewidth}{0.4pt}
\renewcommand{\footrulewidth}{0.4pt}

\begin{document}

% =====================================================================
% NASLOVNICA
% =====================================================================
\begin{titlepage}
    \thispagestyle{empty}
    \centering
    {\fontsize{16pt}{19pt}\selectfont SVEUČILIŠTE U ZAGREBU\\}
    {\fontsize{16pt}{19pt}\selectfont FAKULTET STROJARSTVA I BRODOGRADNJE\\}
    \vfill
    {\fontsize{28pt}{34pt}\selectfont\bfseries {{KOLEGIJ}}\\}
    \vspace{0.5cm}
    {\fontsize{28pt}{34pt}\selectfont\bfseries SEMINAR\\}
    \vfill
    \noindent
    \begin{minipage}[t]{0.48\textwidth}
        \raggedright
        {\fontsize{14pt}{17pt}\selectfont Profesor:\\}
        {\fontsize{14pt}{17pt}\selectfont {{PROFESOR}}\\}
    \end{minipage}
    \hfill
    \begin{minipage}[t]{0.48\textwidth}
        \raggedleft
        {\fontsize{14pt}{17pt}\selectfont\bfseries {{IME_I_PREZIME}}\\}
    \end{minipage}
    \vspace{1cm}
    \begin{center}
        {\fontsize{14pt}{17pt}\selectfont Zagreb, {{GODINA}}.}
    \end{center}
\end{titlepage}

% =====================================================================
% SADRŽAJ I POPISI (rimske stranice)
% =====================================================================
\pagenumbering{Roman}

\tableofcontents
\newpage
\listoffigures
\newpage
\listoftables
\newpage

% =====================================================================
% SADRŽAJ RADA (arapske stranice)
% =====================================================================
\pagenumbering{arabic}

\section{UVOD}
{{UVOD}}

\section{POGLAVLJE}
{{POGLAVLJE_2}}

\section{ZAKLJUČAK}
{{ZAKLJUCAK}}

% =====================================================================
% LITERATURA
% =====================================================================
\newpage
\begin{thebibliography}{99}
\addcontentsline{toc}{section}{LITERATURA}
{{LITERATURA}}
\end{thebibliography}

% =====================================================================
% PRILOZI
% =====================================================================
\newpage
\section*{PRILOZI}
\addcontentsline{toc}{section}{PRILOZI}
{{PRILOZI}}

\end{document}
